\documentclass[10pt]{article}
\usepackage[margin=2.5cm]{geometry}
\usepackage{booktabs}
\usepackage{float}
\usepackage{graphicx}

\title{Labwork 5\\\large HPC -- Gaussian blur convolution}
\author{
Le Nghiem Thanh Thuy - 2540023\\
University of Science and Technology of Hanoi
}
\date{}

\begin{document}
\maketitle

\section{Preparation}
The image is \texttt{astronaut()} from \texttt{skimage.data}, size $512\times512$ with 3
channels. The blur uses the $7\times7$ Gaussian kernel given in the labwork, whose weights
add up to 1003, so after multiplying the window by the kernel I divide the sum by 1003 to
keep the brightness.

Every output pixel is the weighted sum of the $7\times7$ window centred on it. 

I wrote three versions: a CPU one as a reference, a GPU one reading from global memory, and
a GPU one using shared memory. All three give exactly the same image
(\texttt{np.array\_equal} returns \texttt{True}).

\section{Results}
The CPU version takes 8971 ms.

The GPU is kept busy for one second first so that its clock is already up, then every
configuration is launched once as a warm-up and 20 times for the measurement, and I keep
the fastest one.

\begin{table}[H]
\centering
\caption{Kernel time with and without shared memory (Tesla T4 on Google Colab).}
\label{tab:shared}
\begin{tabular}{lrrr}
\toprule
\textbf{Block size} & \textbf{Global memory (ms)} & \textbf{Shared memory (ms)} & \textbf{Speedup} \\
\midrule
$(8, 8)$   & 0.521 & 0.386 & 1.35x \\
$(16, 8)$  & 0.512 & 0.379 & 1.35x \\
$(16, 16)$ & 0.510 & 0.374 & 1.36x \\
$(32, 8)$  & 0.510 & 0.380 & 1.34x \\
$(16, 32)$ & 0.504 & 0.367 & 1.37x \\
$(32, 32)$ & 0.540 & 0.395 & 1.37x \\
\bottomrule
\end{tabular}
\end{table}

\begin{figure}[H]
\centering
\includegraphics[width=\textwidth]{blocksize_shared.png}
\caption{Kernel time without shared memory (blue) and with shared memory (orange).}
\label{fig:shared}
\end{figure}

The best configuration is $(16, 32)$ for both versions: 0.504 ms without shared memory and
0.367 ms with it.

\section{Comment}
Shared memory makes the kernel about 1.35 times faster, and the gain is almost the same for
every block size.

\end{document}
